\documentclass[11pt,a4paper]{article}
\usepackage[utf8]{inputenc} % Not needed with XeLaTeX
\usepackage{fontspec}
\usepackage[margin=0.75in]{geometry}
\usepackage{xcolor}
\definecolor{darknavy}{RGB}{0,0,100}
\usepackage[colorlinks=true, urlcolor=darknavy, linkcolor=darknavy]{hyperref}
\usepackage{enumitem}
% \usepackage{fontawesome}
\usepackage{fontawesome5}
\usepackage{titlesec}
\usepackage{academicons}

% % --- Font: Source Sans Pro via Type1 (works with pdfLaTeX)
% \usepackage[default,scale=1.0]{sourcesanspro} % body + sans by default
% \renewcommand{\familydefault}{\sfdefault}      % make sans the default family
% \usepackage[scaled]{helvet}
\usepackage[default,scale=0.95]{opensans}

% (optional) math that pairs reasonably with sans body
\usepackage{newtxsf} % sans math; or use newtxtext/newtxmath for serif math

\titleformat{\section}{\Large\bfseries}{\thesection}{1em}{}[\titlerule]
\titlespacing{\section}{0pt}{12pt}{6pt}

\usepackage{comicneue}

\begin{document}

\begin{center}
  % {\Huge\textbf{Pavel Popov}}\\[0.5em]
  % {\comicneue\Huge\textbf{Pavel Popov}}\\[0.5em]
  {\fontspec{JetBrains Mono}\Huge Pavel Popov}\\[0.5em]
  {\fontspec{JetBrains Mono}\footnotesize +1 (470) 920-1876 \quad | \quad Center street NW, Atlanta, GA 30318} \quad | \quad {paavali.popov@gmail.com}\\
  {\small \href{https://github.com/paavalipopov}{Github: paavalipopov} \quad | \quad
  % \href{https://www.linkedin.com/in/pavel-popov-83b69a123/}{\faLinkedin\ Pavel Popov} \quad
  \href{https://www.linkedin.com/in/pavel-popov-83b69a123/}{LinkedIn: pavel-popov-83b69a123} \quad | \quad
  \href{https://scholar.google.com/citations?user=4alUD5UAAAAJ}{Google scholar: 4alUD5UAAAAJ}}
\end{center}

\section*{Skills}
\begin{itemize}
[leftmargin=*,noitemsep]
\item \textbf{Programming \& Tools:} Python, PyTorch, Scikit-learn, Pandas, MATLAB, Git, Docker, HPC, SLURM, Spark, SQL, mongoDB, Figma.
\item \textbf{Machine Learning:} Deep learning architecture design, time series analysis, causal inference, explainable AI, distributed data parallel training.
\item \textbf{Data mining:} Frequent pattern analysis, feature engineering, data visualization.
% \item \textbf{Domain Expertise:} Neuroimaging, solid state physics, signal processing.
\end{itemize}

\section*{Open Source Projects}
\noindent\textbf{\href{https://pypi.org/project/ml4fmri/}{ml4fmri}: Python toolkit for benchmarking DL models for fMRI analysis}\\
[-2em]
\begin{itemize}[leftmargin=*,noitemsep]
    \item Unified 13 DL models for multivariate time series analysis under a standardized training API.
    \item Automated an all-model benchmarking pipeline for any input time series for reproducible research.
\end{itemize}

\noindent\textbf{\href{https://pypi.org/project/brainbow/}{brainbow} — Neuroimaging visualization utility}\\
[-2em]
\begin{itemize}[leftmargin=*,noitemsep]
    \item Created a command-line tool for visualizing obscure neuroimaging data formats.
\end{itemize}

\section*{Experience}
\noindent{\textbf{Center for Translational Research in Neuroimaging and Data Science}} \hfill January 2022 --- Now\\
Graduate Research Assistant \hfill \textit{Atlanta, USA}\\
[-2em]
\begin{itemize}[leftmargin=*,noitemsep]
    \item Developed a robust DL classifier for medical diagnostics of functional MRI scans with SOTA performance and 10x speedup to existing benchmarks.
    \item Designed an interpretable AI model to simulate multivariate time-series data, enabling the identification of novel biomarkers and data-driven insights.
    \item Built and optimized frameworks for training and deployment of DL models for neuroimaging analysis in distributed systems.
    \item Coordinated work of cross-functional teams to develop ML tools useful in clinical practice.
    \item Presented and translated technical findings for general audiences, organized and volunteered in social events.
\end{itemize}

\noindent\textbf{Institute of Radioengineering and Electronics} \hfill May 2017 --- September 2020\\
Research Assistant \hfill \textit{Moscow, Russia}\\
[-2em]
\begin{itemize}[leftmargin=*,noitemsep]
    \item Modeled complex signal propagation in novel hardware architectures based on wave logic.
    \item Developed mathematical models for control of high-frequency signal dynamics, providing theoretical validation for prototype hardware controllers.
    \item Ran large-scale numerical simulations to validate theoretical models.
\end{itemize}

\section*{Education}
\noindent\textbf{Georgia State University} \hfill June 2027\\
PhD in Computer Science \hfill
% Machine Learning \& Neuroimaging,
GPA 3.89\\[0.3em]
\textbf{Moscow Institute of Physics and Technology} \hfill June 2020\\
MS in Applied Mathematics and Physics \hfill
% Spintronics,
GPA 3.92\\[0.3em]
\textbf{Moscow Institute of Physics and Technology} \hfill June 2018\\
BS in Applied Mathematics and Physics \hfill
% Solid State Physics,
GPA 3.81

\end{document}
